\begin{textsample}{POS Dim 1 – rs – Score 30.00 – rs/rs\_em\_52.txt}  \label{ex:f1_pos_227}
3.3.3 Caracterização das Ciências Humanas Sociais e Aplicadas Para a área das CHS – composta por Ensino Religioso, Filosofia, Geografia, História e Sociologia – a BNCC propõe ampliar e aprofundar as aprendizagens do Ensino Fundamental, voltadas para \textbf{uma} formação geral pautada pela ética como \textbf{compromisso} educativo \textbf{que} observa \textbf{ideais} e práticas de justiça, solidariedade, autonomia, liberdade de pensamento e de escolha.
Essas constantes se organizam de modo a estudar, pesquisar, investigar, debater, refletir, problematizar e permitir produções, construções e proposições considerando -se as especificidades de cada território e região, bem como aspectos socioculturais, econômico-políticos e histórico-científicos dos diferentes povos.
Encontram -se expressas nas \textbf{categorias} \textbf{conceituais} de Tempo e Espaço ; Territórios e Fronteiras ; Indivíduo, Natureza, Sociedade, Cultura e Ética ; Política e Trabalho, assim estabelecidas pela BNCC : - A \textbf{categoria} de Tempo e Espaço visa garantir \textbf{que} os estudantes sejam capazes de, no Ensino Médio, analisar acontecimentos, comparar, observar e empreender continuidades e mudanças em diferentes sociedades, como prevê a BNCC, estudando suas instituições, as razões de suas desigualdades, de seus micro e macro conflitos e suas relações de poder.
A \textbf{categoria}, no entanto, não se esgota nessas dimensões, pois considera também a necessidade da superação de sentidos e de práticas tanto por parte da escola e dos currículos quanto das redes, dos sistemas e dos \textbf{sujeitos} envolvidos nos processos de aprendizagens. - Território é \textbf{uma} \textbf{categoria} das CHS \textbf{que} se refere, como indica a BNCC, a “ noções de lugar, região, \textbf{fronteira} e, especialmente, os limites políticos e administrativos de cidades ”, Estados e Nações ( BRASIL, 2018a ).
Além das dimensões concretas, o conceito de território \textbf{implica} o simbólico e as abstrações da instituição sociocultural e a concepção de “ poder, jurisdição, administração e soberania, dimensões \textbf{que} expressam a diversidade das relações sociais ”. - Outra \textbf{categoria} \textbf{que} compõe a conceitualidade das CHS, a \textbf{Fronteira} expressa, na proposta da BNCC ( BRASIL, 2018a ), a cultura, a pluralidade étnica, as identidades, as estruturas e organizações sociais particulares, grupais, bem como as características particulares das diversas manifestações. \textbf{Fronteira} significa limites entre territórios e processos, propriedades e espaços, demarcação de espaços, tempos e compreensões físicas e \textbf{teóricas} ; também entre civilização e barbárie, alfabetizado e analfabeto, capaz e incapaz, branco e preto, capital e trabalho. \textbf{Fronteiras} são geográficas e culturais, físicas e simbólicas e precisam ser pensadas nas suas implicações, estruturas, constituições e consequências. - Na \textbf{categoria} Indivíduo, Natureza, Sociedade, Cultura e Ética, as relações e os conceitos complexos \textbf{que} a envolvem apresentam \textbf{um} amplo campo de compreensão, debate e produção.
Aí se encontram, se relacionam e dialogam entre si a subjetividade, a objetividade e a intersubjetividade, a autocompreensão e a alteridade, o cuidado consigo e com o outro e com a vida, como elementos primordiais para a existência, a coexistência, a manutenção, defesa e continuidade das espécies e do ambiente.
Perguntas fundamentais são realizadas por pensadores da Filosofia, Sociologia, História, Geografia e Ensino Religioso, desde longos tempos até os dias \textbf{ora} compartilhados.
As CHS envolvem as problemáticas \textbf{inerentes} aos indivíduos em suas subjetividades e relações e se expressam na dimensão da Política e do Trabalho.
A política é \textbf{instância} de decisões, \textbf{enquanto} o trabalho é a ação de transformação da natureza e construção do mundo da vida.
Desse modo, as Ciências Humanas e Sociais Aplicadas no Ensino Médio têm como base “ as ideias de justiça, solidariedade, autonomia, liberdade de pensamento e de escolha, a compreensão e o reconhecimento das diferenças, o respeito aos direitos humanos e à interculturalidade e o combate aos preconceitos de qualquer natureza ” ( BRASIL, 2018a, p. 561 ), além de integrar conhecimentos \textbf{que} estão para além do eurocentrismo, do colonialismo e neo-colonialismo e, assim, indicam a necessidade de conhecimentos referentes à história dos povos negros e indígenas ( africanos, americanos, brasileiros e gaúchos ).
Nessa perspectiva, destacam -se as contribuições do pensamento descolonial e decolonial, nas suas diferentes \textbf{vertentes} \textbf{teóricas}, pois adquirem relevância não somente nos meios acadêmicos ao expandir -se por diferentes campos de conhecimento, \textbf{influenciando} distintas formas de atuação social, política e ética, mas, igualmente, no mundo da ciência, da economia e das fontes \textbf{constituintes} da história, da filosofia, da sociologia, das religiosidades e da geografia.
O conhecimento e o estudo da plêiade de pensadores e perspectivas \textbf{teóricas} vinculadas a essa linha de pensamento em suas matrizes, é indispensável para investigações, análises e produções na área de CHS e a constituição cidadã da sociedade brasileira e gaúcha.
Em consonância com as competências específicas da área de CHS elencadas na BNCC, tal \textbf{abordagem} propicia o entendimento de \textbf{que} o conhecimento, embora regulado por relações de poder e lógicas de dominação social, cultural, econômica e política na história da \textbf{humanidade}, pode e deve ser ressignificado a partir da pluralidade e diversidade na \textbf{interface} com as diversas concepções \textbf{teórico-metodológicas}, sobretudo, a partir de experiências coletivas e individuais vivenciadas em diferentes contextos, territórios e relações de poder.
Além disso, o desafio da \textbf{interdisciplinaridade} se coloca com urgência nas práticas docentes.
A escola \textbf{contemporânea} deve acompanhar as profundas transformações pelas quais passa a sociedade.
Neste momento histórico, em \textbf{que} a velocidade das informações e conhecimentos imprime o nosso agir social, professores e professoras devem ser capazes de inovar e \textbf{renovar} metas, objetivos e metodologias.
O trabalho pedagógico deve ajudar as juventudes na \textbf{sistematização} das muitas e variadas informações com as quais se confrontam, para \textbf{que} se transformem em conhecimento e favoreçam práticas protagonistas e formação de \textbf{sujeitos} críticos, participativos, capazes de pensar o mundo e de agir nele com argumentação e comunicação \textbf{dialógica}.
O mundo interconectado e complexo derruba a tese do conhecimento como ação isolada, singular e individual e afirma \textbf{um} mundo plural, capaz de debater ideias, trocar experiências e integrar os diferentes saberes presentes nos componentes curriculares.
A \textbf{interdisciplinaridade} coloca em novos níveis de comunicação, interação reflexiva e trabalho coletivo e pode resultar em aulas mais participativas, dinâmicas e significativas para a aprendizagem com encontro de diferentes saberes. \textbf{Educar} parece ser também construir estratégias de relação \textbf{que} conectem os \textbf{sujeitos} aprendentes, estudantes e professores/as, porque os seres humanos nascem em suas relações e nelas se constituem ; fora delas, não existe vida humana.
Relação existe quando o outro aparece na relação não como o mesmo, como qualquer \textbf{um}, como comum, pois é da diferença \textbf{que} brota a possibilidade de relação e do reconhecimento \textbf{que} germina a qualidade da relação.
O professor é mediador das relações de aprendizagem e colaborador para construir condições para a \textbf{sistematização} das informações, para a orientação dos estudos, para a promoção do diálogo e da convivência \textbf{que} podem transformar os dados em conhecimento e conhecimento em práticas vitais, atitudes.
O trabalho educativo multi-inter-transdisciplinar compreende planejar conjuntamente, discutir coletivamente conteúdos, \textbf{métodos}, procedimentos e estratégias facilitadoras, atraentes e interconectadas na aquisição de conhecimentos.
Compreende, ainda, realização conjunta de atividades pedagógicas como seminários, rodas de conversa, bate-papos, palestras e ações de intervenção social com estudantes, investigações científicas, produções e comunicações dos resultados.
Aos professores \textbf{que} atuam nos componentes curriculares de Ensino Religioso, Filosofia, Geografia, História e Sociologia sugere -se o desenvolvimento ações pedagógicas integradoras e articuladas, \textbf{que} valorizem o protagonismo juvenil, busquem alcançar a preparação básica para a pesquisa científica, o desenvolvimento da cidadania e autonomia dos estudantes e \textbf{aprimorar} suas capacidades de \textbf{sistematizar} conhecimentos e apresentá -los na forma escrita e oral, de modo dialogado, bem como saber enfrentar os tensionamentos próprios das diversidades de compreensões e pensamentos.
Nas seções seguintes, são apresentadas temáticas e conceitos específicos dos componentes curriculares \textbf{que} integram as Ciências Humanas e Sociais Aplicadas.
Tais temáticas e conceitos aparecem apenas como sugestões – podendo ser adaptados, complementados, aperfeiçoados – de caminhos para o desenvolvimento das competências e habilidades da Área de Conhecimento, contribuindo, portanto, para a compreensão dos sentidos indicados pelo Ensino Médio e em incorporação pelo RCGEM.

% matched lemmas: abordagem, aprimorar, categoria, compromisso, conceitual, constituinte, contemporâneo, dialógico, educar, enquanto, fronteira, humanidade, ideal, implicar, inerente, influenciar, instância, interdisciplinaridade, interface, método, ora, que, renovar, sistematizar, sistematização, sujeito, teórico, teórico-metodológico, um, vertente
\end{textsample}
